% !PRINTED 615
\weberpage{615}

%<<CONT
oder, wenn man die Zusammensetzung ausführt,
\weq{7}{\varphi(z + 1) \equiv a' \varphi(z) + a\ (\mathrm{mod}\ n)}
für jedes ganzzahlige $z$. Da $\varphi(z)$ den Grad $n$ nicht erreicht, so müssen (\S.\,136) in (7) die Co\"efficienten der einzelnen Potenzen von $z$ auf beiden Seiten congruent sein, und aus der Vergleichung der Co\"efficienten der höchsten Potenzen von $z$ ergiebt sich $a' \equiv 1\ (\mathrm{mod}\ n)$, also
\[
\varphi(z + 1) \equiv \varphi(z) + a\ (\mathrm{mod}\ n).
\]
Setzt man hier $z + 1, z + 2 \ldots z + h - 1$ für $z$, so folgt
\[
\varphi(z + h) = \varphi(z) + ah,
\]
wo $h$ mit jeder beliebigen ganzen Zahl nach dem Modul $n$ congruent sein kann. Darin kann man nun $z = 0$ setzen und erhält, wenn man wieder $z$ für $h$ schreibt und $\varphi(0) = b$ setzt,
\weq{8}{\varphi(z) = az + b.}

Es besteht also die Gruppe $P$ aus lauter linearen Permutationen und unser Satz ist bewiesen.

Wir fügen noch die Bemerkung hinzu, die sich unmittelbar aus dem Anblick der Formeln ergiebt, dass die \gesperrt{cyklische Gruppe ein normaler Theiler einer jeden linearen Gruppe ist, in der sie überhaupt enthalten ist}.

Wenn wir von dem bewiesenen Satze die Anwendung auf die Gruppe der metacyklischen Gleichung machen, so erhalten wir den Satz von \gesperrt{Galois}:

II. \gesperrt{Die Gruppe einer metacyklischen Gleichung von Primzahlgrad ist linear.}

Denn kehren wir zu der Kette der Gruppen (1) zurück, so haben wir gesehen, dass $P_{\mu - 1}$ die cyklische Gruppe ist. Ist $P_{\mu - 1}$ nicht mit $P$ identisch, so ist $P_{\mu - 1}$ ein Normaltheiler von $P_{\mu - 2}$ und also $P_{\mu - 2}$ linear. $P_{\mu - 2}$ ist wieder Normaltheiler von $P_{\mu - 3}$, also ist auch $P_{\mu - 3}$ linear u.\,s.\,f., bis wir zu dem Schluss gelangen, dass auch $P$ selbst linear sein muss. Zugleich ergiebt sich, dass $P_{\mu - 1}$ Normaltheiler aller vorangehenden Gruppen, also auch von $P$ selbst ist.

Es gilt auch der umgekehrte Satz:

III. \gesperrt{Jede irreducible Gleichung von Primzahlgrad, deren Gruppe linear ist, ist metacyklisch.}

Um ihn zu beweisen, genügt es, zu zeigen, dass jede transitive lineare Gruppe $L$ einen Normaltheiler $L'$ hat, dessen Index
%CONT>>
